\documentclass[11pt,a4paper]{article}
\usepackage[utf8]{inputenc}
\usepackage[margin=1in]{geometry}
\usepackage{amsmath,amssymb}
\usepackage{booktabs}
\usepackage{xcolor}
\usepackage{hyperref}
\usepackage{fancyhdr}
\usepackage{enumitem}

\setlength{\headheight}{14pt}

\hypersetup{
    colorlinks=true,
    linkcolor=blue,
    filecolor=magenta,      
    urlcolor=cyan,
    pdftitle={Project Report: LED ON and OFF Circuit Using 555 Timer and BC557},
    pdfpagemode=UseOutlines,
}

\pagestyle{fancy}
\fancyhf{}
\rhead{LED Dimmer / Flasher (555 + BC557)}
\lhead{Technical Project Report}
\rfoot{Page \thepage}

\title{\textbf{\Huge Engineering Project Report}\\[0.5em] \Large LED ON/OFF Fading Dimmer Circuit Using NE555 Timer IC and BC557 PNP Transistor}
\author{\textbf{Bigyanlabs Engineering Team}\\ Consumer Electronics \& Circuit Design Division}
\date{\today}

\begin{document}

\maketitle
\thispagestyle{empty}

\begin{abstract}
This report details the operational principles, component selections, mathematical modeling, and experimental assembly of an LED ON and OFF / Fading Dimmer Circuit. The design utilizes an NE555 precision timer configured as a low-frequency astable multivibrator, generating a continuous pulse train. The output drives a BC557 PNP Bipolar Junction Transistor (BJT) configured as a high-side electronic switch controlling a 6-LED array powered by a $12\,\text{V}$ DC supply. Incorporating an RC smoothing network at the base of the BC557 converts sharp rectangular switching into smooth exponential fading transitions. Complete pinout specifications, component calculation methods, build instructions, and performance verification are included.
\end{abstract}

\newpage
\tableofcontents
\newpage

\section{Introduction \& System Overview}
Light-Emitting Diodes (LEDs) are widely used in low-power lighting, signal indication, and decorative effects. Creating soft pulsing or flashing visual effects enhances user experience without needing programmable microcontrollers.

This project uses an NE555 timer IC to set pulse rates and a BC557 PNP transistor to handle current switching across an array of LEDs.

\section{Bill of Materials}
\begin{table}[h!]
\centering
\caption{Bill of Materials for 555 + BC557 LED Driver Circuit}
\vspace{0.5em}
\begin{tabular}{@{}c l c l p{6cm}@{}}
\toprule
\textbf{S.No.} & \textbf{Component} & \textbf{Qty} & \textbf{Part Number} & \textbf{Circuit Function} \\
\midrule
1 & Timer IC & 1 & NE555N & Astable pulse generator \\
2 & High-Side Transistor & 1 & BC557 (PNP BJT) & High-side current switch for LED array \\
3 & LED Resistors & 6 & $1\,\text{k}\Omega$, 1/4W & Individual current limiters for LEDs \\
4 & Base Limit Resistor & 1 & $56\,\text{k}\Omega$, 1/4W & Base drive current limiter \\
5 & Biasing Resistor & 1 & $47\,\text{k}\Omega$, 1/4W & RC fading network resistor \\
6 & Timing Resistors & 2 & $100\,\text{k}\Omega$, $200\,\text{k}\Omega$ & Astable timing feedback resistors \\
7 & Timing Capacitors & 2 & $10\,\mu\text{F}$, $470\,\mu\text{F}$ & Astable timing ($10\,\mu\text{F}$) and fading ($470\,\mu\text{F}$) \\
8 & Output Indicators & 6 & 5mm LEDs & Visual array output \\
9 & Power Supply & 1 & $12\,\text{V}$ DC Adapter & Main circuit DC supply \\
\bottomrule
\end{tabular}
\end{table}

\section{Transistor Switching Dynamics \& Operating Principles}

\subsection{Astable Multivibrator Oscillation}
The NE555 operates in astable mode. Capacitor $C_1$ ($10\,\mu\text{F}$) charges through $(R_A + R_B) = (200\,\text{k}\Omega + 100\,\text{k}\Omega)$ up to $\frac{2}{3}V_{CC}$ and discharges through $R_B = 100\,\text{k}\Omega$ to Pin 7 down to $\frac{1}{3}V_{CC}$. This outputs a continuous square wave on Pin 3.

\subsection{BC557 PNP High-Side Switching Modes}
A PNP transistor conducts when its Base voltage is driven lower than its Emitter voltage ($V_B < V_E = +12\,\text{V}$):
\begin{enumerate}
    \item \textbf{555 Output Pin 3 is LOW ($0\,\text{V}$):}
    $V_B$ drops below $V_E$, forward-biasing the Base-Emitter junction ($V_{BE} \approx -0.7\,\text{V}$). The BC557 turns **ON** (saturates), passing $+12\,\text{V}$ to the Collector and illuminating the LED array.
    \item \textbf{555 Output Pin 3 is HIGH ($+12\,\text{V}$):}
    $V_B$ equals $V_E$ ($V_{BE} \approx 0\,\text{V}$). The BC557 turns **OFF** (cutoff), cutting off Collector current and extinguishing the LED array.
\end{enumerate}

\subsection{RC Fading / Dimming Transition}
Connecting a $470\,\mu\text{F}$ capacitor in parallel with the $47\text{k}\Omega$ base network creates an RC time constant ($\tau = R \cdot C = 47\,\text{k}\Omega \times 470\,\mu\text{F} \approx 22.09\,\text{seconds}$). As base current charges and discharges exponentially, the LED brightness smoothly fades IN and OUT instead of flashing abruptly.

\section{Assembly \& Construction Guide}
\begin{enumerate}
    \item Place NE555 IC and BC557 transistor on the breadboard.
    \item Connect 555 Pin 8 ($V_{CC}$) and BC557 Emitter to $+12\,\text{V}$ rail. Connect 555 Pin 1 to GND rail.
    \item Connect $200\,\text{k}\Omega$ resistor from $+12\,\text{V}$ to 555 Pin 7. Connect $100\,\text{k}\Omega$ resistor between Pin 7 and Pin 6. Jumper Pin 6 to Pin 2.
    \item Connect $10\,\mu\text{F}$ capacitor from Pin 2/6 to GND.
    \item Connect $56\,\text{k}\Omega$ resistor from 555 Pin 3 to BC557 Base. Connect $470\,\mu\text{F}$ capacitor across the base network to $+12\,\text{V}$.
    \item Connect BC557 Collector to 6 parallel LED branches (each with a $1\,\text{k}\Omega$ resistor in series). Connect LED cathodes to GND.
    \item Power with $12\,\text{V}$ DC to observe smooth LED ON/OFF dimming behavior.
\end{enumerate}

\section{Conclusion}
The NE555 + BC557 LED ON/OFF Dimmer Circuit cleanly illustrates PNP high-side switching and RC exponential fading dynamics, offering a simple analog visual effect generator.

\end{document}
